\section{A first choice and a place are different outcomes}\label{sec:ranks}
Was the lower placement count the price of better assignments for other applicants? To find out, we ranked each assignment on the person's original list, before either shortlist was made. A renter's order is modeled. For a pupil we use the rank the family entered, distinguishing the ministry's appended continuity option. An unplaced person stays in the denominator at every cutoff. Conditioning on being matched would omit the very people whose losses we want to measure.

\begin{table}[ht]\centering\small
\caption{Change in assignments at or above each rank. NYC compares Haiku with uniform menus, in geographically weighted percentage points. Chile compares the rank-informed predictor with the rank-blind predictor, in pupils. The first three school columns count only family-entered choices; the last includes any assignment, including an administrative fallback.}\label{tab:ranks}
\begin{tabular}{@{}lrrrr@{}}\toprule
 & First choice & Top three & Top five & Any placement\\\midrule
NYC ordinary (pp) & $+2.00$ & $+4.22$ & $+4.99$ & $-2.71$\\
NYC enriched (pp) & $+2.87$ & $+6.69$ & $+7.72$ & $-5.40$\\\addlinespace
Chile 2022 (pupils) & $+385$ & $+70$ & $-184$ & $-468$\\
Chile 2023 & $+515$ & $+116$ & $-181$ & $-517$\\
Chile 2024 & $+747$ & $+14$ & $-303$ & $-481$\\
Chile 2025 & $+571$ & $-38$ & $-254$ & $-321$\\\bottomrule
\end{tabular}
\end{table}

In every panel the new policy produces more first choices and fewer placements. Neither policy therefore dominates the other at every rank cutoff, in the sense of \citet{featherstone}. This conclusion also holds if we include each administrative fallback at its final list position: the first-choice and any-placement differences are unchanged. In the ordinary and enriched NYC markets, the paired 95\% simulation intervals for the first-choice gains are $[1.67,2.33]$ and $[2.41,3.32]$ percentage points. These are ranks on the modeled full order, not evidence that particular renters would have signed a lease.

There is precedent for this pattern under a different policy. In Chilean simulations \citet[Table~6]{correaschool} found that changing from course-specific lotteries to a common lottery increased both first choices and nonassignment. \citet{trace2026} also compare tie rules using several welfare measures. Our counterfactual keeps the priority order and quota fixed, changes the predicted shortlist, and observes which affected families had the same list under both policies.

A simple accounting exercise puts a price on the disagreement without estimating anyone's utility. Give one point for any placement and an additional $\lambda$ points for a first choice. On the four school cohorts together, the change in this index is $-1{,}787+2{,}218\lambda$. Someone using this objective would prefer the rank-informed rule if and only if $\lambda>1{,}787/2{,}218\simeq0.806$, with a tie at equality. Nothing in the data identifies that bonus or implies it is the same for all families. Counting assignments alone implicitly sets it to zero.

Who moved up? Ordering the outside option last and an appended continuity option after all family-entered choices, 8,076 Chilean applicants with edited lists got a strictly better result and 2,200 got a worse one. Among applicants with identical lists, 2,254 improved and 7,215 worsened. Each of the four years has this same split in direction. These are paired changes between two full-market replays. The two groups differ in list length and exposure to competition. Comparing their rates would not identify a treatment effect of editing one's own list.
